% The research corpus (ci/research.json): a mathematics monograph in
% amsbook (TeX Live's amsbook-template.tex does not build as shipped: its
% \include{} placeholders name no file).
\documentclass{amsbook}
\usepackage{amssymb}
\newtheorem{theorem}{Theorem}[chapter]
\newtheorem{lemma}[theorem]{Lemma}
\theoremstyle{definition}
\newtheorem{definition}[theorem]{Definition}
\theoremstyle{remark}
\newtheorem{remark}[theorem]{Remark}
\numberwithin{section}{chapter}
\numberwithin{equation}{chapter}
\makeindex
\begin{document}
\frontmatter
\title{Lectures on Finite Groups}
\author{L. Algebraist}
\address{Department of Mathematics, University of Examples}
\email{algebraist@example.org}
\subjclass[2020]{Primary 20D05; Secondary 20C15}
\keywords{finite groups, Sylow theorems, representations}
\date{January 1, 2026}
\begin{abstract}
A short monograph in the American Mathematical Society's book class.
\end{abstract}
\maketitle
\tableofcontents
\chapter*{Preface}
These notes grew out of a one-semester course.
\mainmatter
\chapter{Groups}
\section{Definitions}
\begin{definition}\label{def:group}
A \emph{group} is a set $G$ with an associative operation, an identity
$e$ and inverses: $g g^{-1} = g^{-1} g = e$ for all $g \in G$.
\end{definition}
\begin{theorem}[Lagrange]\label{thm:lagrange}
If $H \le G$ and $G$ is finite, then $|H|$ divides $|G|$:
\begin{equation}\label{eq:lagrange}
  |G| = [G:H]\,|H| .
\end{equation}
\end{theorem}
\begin{proof}
The left cosets of $H$ partition $G$ and each has $|H|$ elements.
\end{proof}
\section{Sylow's theorems}
\begin{theorem}[Sylow]
Let $|G| = p^a m$ with $p \nmid m$. Then $G$ has a subgroup of order
$p^a$, any two are conjugate, and their number $n_p$ satisfies
$n_p \equiv 1 \pmod p$ and $n_p \mid m$.
\end{theorem}
\begin{remark}
Theorem~\ref{thm:lagrange} and \eqref{eq:lagrange} are used throughout.
\end{remark}
\chapter{Representations}
\section{Characters}
For a representation $\rho\colon G \to \mathrm{GL}(V)$ the character is
$\chi(g) = \operatorname{tr}\rho(g)$, and
\[
  \langle \chi, \psi \rangle = \frac{1}{|G|} \sum_{g\in G} \chi(g)\overline{\psi(g)} .
\]
\begin{lemma}[Schur]
A $G$-map between irreducible representations is zero or an isomorphism.
\end{lemma}
\begin{proof}
Its kernel and image are subrepresentations.
\end{proof}
\appendix
\chapter{Notation}
$\mathbb{Z}$, $\mathbb{Q}$, $\mathbb{C}$: the integers, rationals and
complex numbers; $\mathfrak{S}_n$ the symmetric group.
\backmatter
\begin{thebibliography}{9}
\bibitem{serre} J.-P. Serre, \emph{Linear Representations of Finite
Groups}, Springer, 1977.
\bibitem{isaacs} I. M. Isaacs, \emph{Character Theory of Finite Groups},
Academic Press, 1976.
\end{thebibliography}
\end{document}
